\documentclass[11pt]{article}
\usepackage{icmlko}
\begin{document}

\title{계절은 1등이 아니라 늘 있었다 --- 쪼갬 몫으로 중요도를 읽으면
안 된다}
\author{월드모델 실험실}
\date{노트 479 · 2026년 8월 3일}
\maketitle

\begin{abstract}
노트 478 이 덤으로 하나 적었다 --- \emph{``계절이 제일 많이 쪼개지는
축이다($0.0907$). 노트 447 이 넣은 축인데 두 해 묵은 축들보다 더
쓰인다.''} 그런데 계절이 판에 준 이득은 $+0.0041$ 뿐이다. \textbf{1등
축치고 작다.} 관찰 자체를 의심했다.

의심한 것은 \emph{값 종류 수}였다 --- 계절은 $0\sim1$ 연속이니 언제든
자를 데가 있다는 것. \textbf{틀렸다}(스피어만 $+0.259$, p $0.12$).

\textbf{덮음이었다.}

\begin{center}
\begin{tabular}{lr}
\hline
쪼갬 몫과의 스피어만 & \\
\hline
덮음 & $\mathbf{+0.761}$ \\
값 종류 수 & $+0.259$ \\
\hline
덮음 $|$ 값 종류 통제 & $\mathbf{+0.785}$ \\
값 종류 $|$ 덮음 통제 & $+0.386$ \\
\hline
\end{tabular}
\end{center}

\texttt{goods\_scale} 은 값이 $3{,}319$종으로 제일 많은데 덮음 $62\%$ 라
몫이 $0.0566$ 이고, 계절은 $1{,}485$종인데 \textbf{덮음 $100\%$} 라
$0.0910$ 으로 1등이다. \textbf{계절이 1등인 것은 유용해서가 아니라
언제나 있어서다.} 나무는 축이 없는 행에서는 그 축으로 못 자른다.
\end{abstract}

\begin{figure}[t]
\centering
\includegraphics[width=\textwidth]{figs/f1.pdf}
\caption{\textbf{가} --- 쪼갬 몫이 덮음을 따라간다. 값 종류가 많은
파란 점들이 덮음이 낮으면 아래에 깔린다. \textbf{나} --- 덮음으로
나누면 순위가 뒤집히는데, 새 1등들이 덮음 $7\sim26\%$ 라 이번엔 반대로
과교정이다.}
\end{figure}

\section{나눠도 안 된다}

그러면 덮음으로 나누면 되나 --- \emph{``있을 때 얼마나 쪼개나''}. 계절이
$1$위에서 \textbf{$18$위}로 떨어진다. 그런데 새 상위권이
\texttt{fund\_cat}($13\%$) · \texttt{cal\_holiday\_gap}($26\%$) ·
\texttt{wiki\_level}($7\%$) 로 \textbf{전부 저덮음 축}이다. 드물게
있는 축은 그 좁은 무리 안에서 여러 번 쪼개지므로 나눈 값이 부풀려진다.

\textbf{날것은 덮음에 쏠리고 나눈 것은 반대로 쏠린다. 어느 쪽으로도
중요도 순위가 안 나온다.}

\section*{조항 --- 구조 자의 쓸 자리를 좁힌다}

노트 477 이 구조 자(잎 수 · 쪼갬 몫 · 쓰인 축)를 만들었고 노트 478 이
그것으로 lr 물음을 풀었다. \textbf{그때는 옳게 썼다} --- 같은 축 집합을
설정끼리 견줬으니 덮음이 양쪽에서 같았다.

\begin{quote}
\textbf{조항} --- 구조 자는 \textbf{같은 축 집합을 설정끼리 견주는 데만}
쓴다. \textbf{축끼리 견주는 데는 안 쓴다} --- 쪼갬 몫은 덮음이 지배하고
(부분 rho $+0.785$), 덮음으로 나누면 반대로 뒤집힌다.
\end{quote}

\section{노트 478 의 덤 한 줄을 물린다}

\emph{``계절이 두 해 묵은 축들보다 더 쓰인다''}는 문장을 물린다. 정확히는
\textbf{``계절은 덮음이 $100\%$ 인 몇 안 되는 축이다''}이다 --- 값 열
$37$개 중 덮음 $95\%$ 를 넘는 것은 계절 · \texttt{target\_breadth} ·
\texttt{venue\_prominence} 셋뿐이다.

그러면 계절의 판 이득이 $+0.0041$ 로 작은 것도 이상하지 않다.
\textbf{많이 쪼개진다고 많이 버는 게 아니다} --- 늘 있는 축은 늘
후보에 오를 뿐이다.

\section*{남은 것}

\begin{itemize}
\item \textbf{덮음을 안 타는 중요도 자} --- 결국 \emph{빼고 다시 재기}
말고는 없나. 축 $37$개를 하나씩 빼면 관문이 $37$번인데, 노트 475 의
값싼 앞자를 여기에도 쓸 수 있나.
\item \textbf{덮음 $95\%$ 넘는 축이 셋뿐이다} --- 나머지 $34$개는
행의 절반 이상에서 없다. 그게 이 자료의 제일 큰 구조적 사실일 수 있다.
\item \textbf{계절의 $+0.0041$ 은 어디서 오나} --- 노트 447 은 애니에서
왔다고 적었다($+0.0264 \cdot 12/12$). 쪼갬 몫과 무관하다는 것만 새로
안다.
\end{itemize}

\end{document}
